% Shared section: interpretation and limitations of the region vs. income analysis
Several caveats apply to the comparison between income and region. First, region is not a mechanism: it bundles culture, history, institutions, social ties, and possibly cross-cultural differences in how people use response scales \citep{chen_response_1995,diener_personality_2003,exton_comparing_2015}. If regional differences partly reflect response styles, then region predicts \textit{measured} well-being without necessarily predicting \textit{experienced} well-being. Existing evidence suggests that self-reported well-being carries information about true well-being \citep{sandvik_subjective_1993,kaiser_scientific_2022}, and studies exploiting migrants suggest that national differences partly reflect a cultural component of well-being rather than mere translation artifacts \citep{senik_french_2014,exton_comparing_2015}, but the present analysis cannot separate these interpretations. Anchoring vignettes \citep{kapteyn_life_2010} or measures of emotional well-being would help.

Second, the advantage of region rests largely on Latin America and the former Eastern Bloc. These are well-documented deviations from the income gradient, but they are not independent of the classification: grouping countries into continents or World Bank regions, which do not isolate these groups, weakens the result. The fairest summary is that GDP per capita leaves much of the cross-country variation unexplained, and that a few large regional deviations explain more of it than income does in the WVS.

Third, the analysis pools country-years that are not independent, and $R^2$ decompositions carry no information on causality. The within-country evidence shows that growth is associated with rising life satisfaction, so the weak cross-sectional relation does not imply that income is irrelevant for well-being.

Fourth, and most importantly, the WVS and Gallup disagree on the strength of the income gradient. In Gallup, income predicts life evaluations about as well as region, or better.
